problem:  
Let \( (M=G/H,g) \) be a compact, simply connected **naturally reductive** Riemannian manifold with reductive decomposition \( \mathfrak{g}=\mathfrak{h}\oplus\mathfrak{m}\) and an \(\mathrm{Ad}(H)\)-invariant inner product on \(\mathfrak{m}\).  
Denote by \( \tilde{\nabla} \) the **canonical connection** with torsion \(T(X,Y)=-[X,Y]_{\mathfrak m}\), and by \( \mathfrak{hol}(\tilde{\nabla}) \subset \mathfrak{so}(\mathfrak m) \) its canonical holonomy algebra.  
Let \( \rho: H \to \mathrm{SO}(\mathfrak m)\) be the **isotropy representation**, with image algebra \( \rho_*(\mathfrak h) \subset \mathfrak{so}(\mathfrak m)\).  

Define the **canonical holonomy generating algebra**
\[
\mathfrak{h}_c := \mathrm{Lie} \langle R^{\tilde{\nabla}}(X,Y) : X,Y\in \mathfrak{m}\rangle,
\]
where \(\langle \cdot \rangle\) denotes the smallest Lie algebra containing all canonical curvature endomorphisms.  By Ambrose–Singer, \(\mathfrak{h}_c = \mathfrak{hol}(\tilde{\nabla})\), and for naturally reductive spaces this algebra is \(\tilde{\nabla}\)-parallel.

**Refined Open Problem (Canonical Holonomy–Isotropy Rigidity Conjecture):**  
Assume that the canonical holonomy satisfies  
\[
\mathfrak{hol}(\tilde{\nabla}) = \mathfrak{so}(\mathfrak{m}),
\]
that is, the canonical connection has **maximal possible holonomy** acting irreducibly on \(\mathfrak{m}\).  
Must \( (M,g) \) then be **normal homogeneous**, i.e. its naturally reductive metric arises from a bi-invariant metric on a larger transitive Lie group?

Equivalently, does maximal canonical holonomy imply that the torsion \(T\) integrates in a symmetric way, forcing \( (M,g) \) to be normal homogeneous?

---

Why is it a “good” problem:  

1. **Conceptually clear yet strikingly refined:**  
   The condition “\(\mathfrak{hol}(\tilde{\nabla})=\mathfrak{so}(\mathfrak{m})\)” is both transparent and extreme—it states that the canonical connection’s curvature spans the entire space of skew-symmetric endomorphisms. Whether such *maximal canonical holonomy* geometrically forces normal homogeneity is entirely unknown and probes the upper holonomy boundary of naturally reductive spaces.

2. **Bridges holonomy theory, torsion geometry, and homogeneous structures:**  
   The problem ties together three major ideas:  
   (a) **Canonical holonomy** from connection theory,  
   (b) **Naturally reductive torsion** from homogeneous geometry, and  
   (c) **Normal homogeneous restriction** from Lie-theoretic metric construction.  
   Progress requires understanding subtle curvature–torsion interactions and their representation effects on \(\mathfrak{so}(\mathfrak{m})\).

3. **Positions itself between known extremes:**  
   - For **bi-invariant metrics**, \(\mathfrak{hol}(\tilde{\nabla})\) is small (equal to the isotropy algebra).  
   - For some special **non-normal** naturally reductive metrics, it can be large but not maximal.  
   Determining whether *maximal* canonical holonomy can occur without normal homogeneity functions as a precise rigidity threshold problem within homogeneous Riemannian geometry.

4. **Potential for major geometric insight:**  
   If true, the conjecture would reveal a new holonomy-type characterization of normal homogeneous manifolds, providing a bridge between the algebraic completeness of \(\mathfrak{so}(\mathfrak{m})\) and the geometric symmetry of normality.  
   If false, any counterexample would expose new families of highly curved, torsion-rich naturally reductive spaces with unexpectedly large canonical holonomy—objects previously unknown.

5. **Simple formulation, great depth:**  
   The statement involves only standard notions—canonical holonomy, naturally reductive decomposition, normal homogeneity, and the symbol \(\mathfrak{so}(\mathfrak{m})\)—yet deciding its truth demands profound understanding of the interplay between algebraic holonomy representations and global group actions.  
   This blend of **extreme simplicity in form** and **depth of structural implication** exemplifies the essence of a “good” research problem.